\chapter*{Introducción}
\addcontentsline{toc}{chapter}{Introducción}
\markboth{Introducción}{Introducción}

El Método de los Elementos Finitos (MEF) constituye una de las herramientas centrales del análisis estructural y de la ingeniería computacional \autocite{zienkiewicz2013fem,cook2002concepts}, y su enseñanza forma parte del núcleo curricular de las carreras de ingeniería civil y mecánica \autocite[p.~1160]{perezsantiago2023fem}. Su comprensión exige articular tres planos que suelen presentarse de forma fragmentada al estudiante: la formulación matemática continua (elasticidad lineal, ecuaciones de equilibrio), su discretización mediante elementos isoparamétricos y cuadratura numérica, y la interpretación física de los resultados (campos de desplazamiento, deformación y tensión). EduFEM se propone como un software educativo de escritorio que integra estos tres planos en un único entorno interactivo, orientado a problemas bidimensionales de elasticidad lineal estática con elementos cuadriláteros Q4 y Q9.

\section*{Planteamiento del problema}

La enseñanza del MEF enfrenta una tensión recurrente entre el rigor formal y la intuición operativa. Por un lado, los textos clásicos desarrollan la teoría con un grado de abstracción considerable, donde las funciones de forma, la matriz Jacobiana, la matriz de deformación-desplazamiento y la integración de Gauss aparecen como pasos de una sucesión algebraica cuya conexión con la geometría concreta de un elemento no siempre resulta evidente para quien recién se inicia \autocite{bathe2014fem,reddy2006introduction}. Por otro lado, el software comercial de análisis por elementos finitos opera, desde la perspectiva del estudiante, como una caja negra: el usuario define geometría, material y cargas, y obtiene resultados, pero el procesamiento intermedio que transforma esas entradas en tensiones permanece oculto \autocite[p.~1160]{perezsantiago2023fem}. El antecedente educativo más directo de este trabajo lo describía ya en 1998 en los mismos términos: con los programas comerciales el estudiante queda limitado a introducir los datos del problema e interpretar los resultados, sin acceso a los detalles del análisis \autocite[p.~246]{suarez1998edelas2d}.

Esta opacidad tiene una consecuencia pedagógica concreta: el estudiante puede aprender a operar una herramienta profesional sin formar criterio sobre lo que obtiene, es decir, sin saber si una deformada tiene sentido, si la malla es admisible, si las tensiones nodales cierran con las de los puntos de Gauss, si las reacciones equilibran la carga o si el elemento bloquea. Un estudio con expertos de la industria y de la academia confirma que ese criterio es lo que la formación tiende a dejar de lado —privilegia las destrezas de operación sobre los conceptos fundamentales— y reclama que el egresado sepa planificar, verificar y validar sus análisis, no solo ejecutarlos \autocites[p.~1159]{perezsantiago2023fem}[p.~1171]{perezsantiago2023fem}; el mismo estudio recoge una encuesta profesional según la cual solo el $31\,\%$ de los ingenieros se considera conocedor del análisis por elementos finitos \autocite[p.~1160]{perezsantiago2023fem}. A esto se suma que las herramientas comerciales suelen ser costosas, no estar diseñadas con un propósito didáctico explícito \autocite[p.~1160]{perezsantiago2023fem} y ofrecer una localización al español limitada como recurso de aula. La brecha, en consecuencia, es doble: existe una dificultad de aprendizaje genuina vinculada a la naturaleza abstracta del método, y una insuficiencia de herramientas que la atiendan exponiendo el proceso de cálculo y la respuesta de manera transparente y guiada. Esa es la \emph{situación problemática} de la que parte el trabajo: en la investigación tecnológica el problema no se elige libremente, sino que resulta de interpretar y diagnosticar una realidad concreta \autocite[p.~83]{garciacordoba2005tecnologica}.

En términos formales, el \textbf{problema científico} de esta investigación se formula como la siguiente interrogante: ¿de qué manera el uso del software de análisis estructural como caja negra podrá afectar en el bajo criterio para interpretar la respuesta estructural obtenida por el MEF en los estudiantes de la Carrera de Ingeniería Civil de la Universidad Autónoma Tomás Frías en la gestión 2026? Es la única formulación del problema que el documento emplea: la \autoref{sec:matriz-consistencia} la retoma textualmente y la \autoref{sec:validacion-hipotesis} responde la parte que este trabajo contrasta. La variable causa es la forma en que el software expone el canal de cálculo y la respuesta estructural —como caja negra, o de manera transparente, interactiva y numéricamente verificable—; la variable efecto es el criterio del estudiante para interpretar y juzgar esa respuesta: desplazamientos, tensiones, reacciones, calidad de la malla y fenómenos numéricos. La insuficiencia de una herramienta que exponga el canal con esos atributos es la contradicción que origina el trabajo \autocite[p.~7]{alvarez_metodologia}. El problema se delimita en lo disciplinar a la elasticidad lineal plana con elementos cuadriláteros Q4 y Q9 (véase \emph{Alcance y limitaciones}); en lo institucional, a la formación de grado en ingeniería civil, con la Carrera de Ingeniería Civil de la Universidad Autónoma Tomás Frías como destinataria inmediata de la herramienta; en lo espacial, al aula y al equipo personal del estudiante, sin requerir laboratorio de cómputo ni licencias; y en lo temporal, al desarrollo, verificación y validación del software realizados durante la gestión 2026.

El \textbf{objeto de estudio} \autocite[p.~8]{alvarez_metodologia} es el proceso de enseñanza y aprendizaje del análisis estructural de medios continuos en elasticidad lineal plana por el Método de los Elementos Finitos —el canal de cálculo que va del mapeo isoparamétrico a la recuperación de tensiones con elementos cuadriláteros Q4 y Q9— en la formación de grado en ingeniería civil. El \textbf{campo de acción}, la parte de ese objeto que el trabajo se propone mejorar \autocite[p.~13]{alvarez_metodologia}, son los medios didácticos con que el estudiante aprende a interpretar y verificar la respuesta estructural: el software educativo que expone el canal de cálculo sobre el modelo del estudiante y la memoria de cálculo que documenta el resultado.

\section*{Justificación}

El desarrollo de EduFEM se justifica en varios planos complementarios. En el plano académico, la herramienta busca cerrar la distancia entre la formulación teórica del MEF y su comportamiento numérico, haciendo visibles los pasos intermedios que el software profesional encapsula. La literatura sobre enseñanza del método reporta que la simulación interactiva y la manipulación directa de los objetos matemáticos del MEF conducen a una comprensión más intuitiva de sus procedimientos \autocites[p.~157]{lee2015interactive}[pp.~872, 874]{lee2015eigenmodes}, que construir paso a paso las herramientas de cálculo junto con el alumno resulta más eficaz que la sola exposición teórica \autocite[p.~17]{bishay2020teaching}, y la revisión de una década de estudios cuasiexperimentales encuentra evidencia sólida de que las simulaciones por computadora mejoran la instrucción tradicional, con un efecto que depende de cómo se visualiza la información y del apoyo instruccional que la acompaña \autocite[p.~136]{rutten2012learning}. Exponer el mapeo isoparamétrico, el Jacobiano, la matriz de deformación y la cuadratura de Gauss sobre el elemento real que el estudiante construyó, y mostrar la respuesta que producen, permite anclar conceptos abstractos en una geometría tangible y juzgar el resultado con conocimiento de cómo se obtuvo.

En el plano práctico, EduFEM está concebido como recurso de apoyo a la cátedra: es gratuito y de código abierto —su código fuente está disponible públicamente en línea\footnote{Repositorio del proyecto: \url{https://github.com/Sucullani/SoftwareED}}—, está íntegramente en español y organiza su interfaz según las tres fases canónicas del análisis (pre-proceso, proceso y post-proceso), lo que reproduce el flujo de trabajo profesional al tiempo que lo hace pedagógicamente legible. La generación automática de una memoria de cálculo —un documento que reconstruye el procedimiento completo, con sustitución numérica paso a paso, sobre los valores del modelo del propio usuario— ofrece un puente entre la exploración interactiva y el cálculo verificable a mano. Esta capacidad —que no se documenta entre los antecedentes de software educativo revisados (\autoref{sec:antecedentes})— convierte a la herramienta en un recurso de verificación: el estudiante puede reproducir cada operación con lápiz y papel o en una hoja de cálculo y contrastar sus resultados con los del programa, elemento por elemento.

En el plano metodológico, la decisión de construir un software propio en lugar de adoptar uno existente responde a la necesidad de controlar la transparencia del motor de cálculo y de alinear cada componente de la interfaz con un objetivo didáctico. Un motor desarrollado con bibliotecas numéricas de código abierto \autocite{harris2020numpy,virtanen2020scipy} permite exhibir las matrices y operaciones intermedias sin las restricciones de un producto cerrado, y verificar su precisión frente a referencias reconocidas.

\section*{Objetivos}

El objetivo general de este trabajo es desarrollar, en el lenguaje de programación Python y sobre bibliotecas numéricas de código abierto, EduFEM, un software educativo de escritorio para el análisis por el Método de los Elementos Finitos (MEF) de problemas bidimensionales de elasticidad lineal (tensión y deformación plana) con elementos cuadriláteros isoparamétricos Q4 y Q9, que integre la formulación matemática, la construcción y visualización interactiva del modelo y la verificación numérica, para mejorar el criterio de los estudiantes de ingeniería civil en la interpretación de la respuesta estructural obtenida por el MEF.

Los objetivos específicos que concretan este propósito son los siguientes:

\begin{enumerate}
    \item Diagnosticar, a partir de la literatura y de los antecedentes de software educativo, la brecha entre operar un software de análisis y saber juzgar su respuesta, para delimitar los atributos que la herramienta debe tener.
    \item Fundamentar teóricamente el MEF para la elasticidad lineal plana con elementos isoparamétricos Q4 y Q9 —formulación variacional, mapeo isoparamétrico, integración numérica, ensamblaje, recuperación de tensiones, calidad de malla, fenómenos numéricos y criterios de verificación y validación— y los principios de aprendizaje que orientan el diseño, como base común del motor de cálculo, de los módulos educativos y de la memoria de cálculo.
    \item Desarrollar el motor de cálculo MEF 2D con elementos Q4 y Q9, la interfaz gráfica organizada en pre-proceso, proceso y post-proceso sobre un único lienzo compartido, los módulos educativos que exponen paso a paso las etapas del canal de cálculo sobre el modelo real, y la memoria de cálculo automática en PDF con interoperabilidad de datos (DXF, CSV), para exponer cada etapa del canal y cada resultado sobre el modelo del estudiante.
    \item Verificar y validar el motor de cálculo mediante el método de soluciones manufacturadas (MMS) y casos de referencia clásicos (viga de Timoshenko, membrana de Cook), contrastando contra soluciones analíticas y un software comercial (SAP2000), para que la respuesta que el software expone sea correcta.
    \item Validar el contenido del software frente a las destrezas de interpretación y verificación de la respuesta que los expertos exigen al egresado, mediante una tabla de especificaciones, y su coherencia con los principios de aprendizaje que fundamentan el diseño, para establecer su potencial para formar criterio.
\end{enumerate}

\section*{Hipótesis y preguntas de investigación}

La \emph{hipótesis} de este trabajo es que, bajo el desarrollo de EduFEM, la exposición transparente e interactiva del canal de cálculo y de la respuesta estructural permitirá mejorar el criterio para interpretar la respuesta obtenida por el MEF en los estudiantes de la Carrera de Ingeniería Civil de la Universidad Autónoma Tomás Frías. En la nomenclatura de tres variables, la variable independiente es la forma en que se expone el canal de cálculo y la respuesta —como caja negra o de manera transparente, interactiva y verificable—, la variable dependiente es el criterio del estudiante para interpretar la respuesta, y la variable interviniente es EduFEM, la solución tentativa con que se manipula la primera. Por tratarse de un trabajo de desarrollo tecnológico, la hipótesis es la solución tentativa a un problema concreto y su criterio de veracidad es la efectividad en la práctica \autocite[p.~85]{garciacordoba2005tecnologica}, y se comprueba por etapas. En esta etapa se comprueba en su parte de diseño y contenido, mediante tres cláusulas con criterios fijados a priori en la \autoref{sec:validacion-resultados}: (a) \emph{exposición}: el software expone cada etapa del canal de cálculo en un módulo educativo o en la memoria, y la respuesta en el post-proceso, sobre las tres fases del análisis en un mismo lienzo; (b) \emph{corrección}: su motor reproduce las tasas teóricas de convergencia del desplazamiento y acota el error por debajo del $3\,\%$ en la flecha y del $1\,\%$ en la tensión normal frente a la solución analítica de la viga de Timoshenko y frente al modelo de SAP2000, y evidencia sobre la membrana de Cook el bloqueo por cortante del Q4 frente a la convergencia del Q9; y (c) \emph{validez de contenido y coherencia}: cada destreza de interpretación y verificación de la respuesta que los expertos exigen tiene un instrumento en el software, y cada principio de aprendizaje que fundamenta el diseño tiene una encarnación comprobable. Cada cláusula puede resultar falsa —una etapa sin exposición, un motor con un error de orden, una destreza sin instrumento o un principio sin encarnación la refutarían—. Que esa exposición mejore de hecho el criterio del estudiante es la parte de la hipótesis que esta etapa no contrasta: se fundamenta en la literatura que reporta ese efecto con herramientas de la misma familia (\autoref{sec:fundamentos-pedagogicos}) y su prueba de campo se deja diseñada en las recomendaciones.

De esta hipótesis se desprenden las siguientes preguntas de investigación, una por cláusula, que el trabajo procura responder:

\begin{itemize}
    \item \textbf{PI-1} (exposición). ¿Qué organización de la interfaz, de los módulos educativos y de la memoria de cálculo expone el canal de cálculo del MEF y la respuesta estructural de manera transparente e interactiva sobre el modelo del estudiante, en lugar de como caja negra?
    \item \textbf{PI-2} (corrección). ¿Es posible construir un motor de cálculo MEF 2D con elementos Q4 y Q9 cuya precisión numérica sea verificable frente a soluciones analíticas y software comercial de referencia, y que reproduzca los fenómenos numéricos que la teoría predice, como el bloqueo por cortante del Q4 frente a la convergencia del Q9?
    \item \textbf{PI-3} (validez de contenido). ¿Cubre el software las destrezas de interpretación y verificación de la respuesta que los expertos exigen al egresado, y encarna los principios de aprendizaje que fundamentan su diseño?
\end{itemize}

\section*{Metodología de la investigación}

Por su finalidad, este trabajo es de carácter \emph{propositivo}: su núcleo es la propuesta y construcción de un artefacto (EduFEM) que resuelve un problema formativo concreto. Por su naturaleza se inscribe, además, en la \emph{investigación aplicada} de tipo \emph{tecnológico}: su producto es un artefacto de software cuya construcción y evaluación constituyen el núcleo del estudio. Por el abordaje del problema, el enfoque es \emph{cuantitativo} en la verificación y validación del motor, pues la corrección del artefacto se establece midiendo tensiones, desplazamientos, errores numéricos y tasas de convergencia frente a referencias de exactitud conocida, y \emph{documental} en la validación de contenido, que confronta los instrumentos del software con las destrezas que un consenso publicado de expertos exige y con los principios de aprendizaje de la literatura \autocite[p.~168]{lee2015interactive}. Según su nivel u objetivo, la investigación es \emph{descriptiva} y \emph{comparativa}: caracteriza el comportamiento numérico del modelo, contrasta de forma sistemática el desempeño de los elementos Q4 y Q9 y describe el grado en que el artefacto expone lo que la hipótesis postula; el nivel explicativo —el efecto de esa exposición sobre el criterio del estudiante— corresponde a la prueba de campo que se deja diseñada.

El diseño metodológico —el modelo de simulación numérica que sostiene la verificación, las variables y su operacionalización, la matriz de consistencia, la población, los casos de estudio y su criterio de selección, el procedimiento con sus instrumentos numéricos y documentales, y los criterios de aceptación con que se contrasta la hipótesis— se desarrolla en detalle en la \autoref{sec:metodologia}.

\section*{Alcance y limitaciones}

El alcance de EduFEM está delimitado de manera explícita para garantizar profundidad sobre un dominio acotado. En este documento, la expresión \emph{análisis estructural} se emplea en su acepción restringida de análisis tenso-deformacional de elementos estructurales modelados como medios continuos bidimensionales mediante la teoría de la elasticidad lineal; no comprende el análisis de sistemas estructurales discretos —pórticos o reticulados—, ni la formulación de placas, cáscaras o problemas tridimensionales. La herramienta resuelve problemas \emph{bidimensionales} de \emph{elasticidad lineal estática}, en las dos idealizaciones clásicas de la mecánica de sólidos: \emph{tensión plana} (apropiada para cuerpos delgados cargados en su plano) y \emph{deformación plana} (apropiada para cuerpos prismáticos largos con sección y carga invariantes a lo largo del eje longitudinal) \autocite[p.~11]{timoshenko1970elasticity}. La discretización se realiza con \emph{elementos cuadriláteros isoparamétricos} de cuatro nodos (Q4) y de nueve nodos (Q9), que cubren respectivamente la interpolación bilineal y la cuadrática completa, y permiten ilustrar de manera contrastada los efectos del orden de interpolación sobre la precisión \autocites[pp.~137-138]{hughes2000fem}[pp.~164, 183-184]{onate2009structural}. El producto es un \emph{software de escritorio} de carácter \emph{educativo}, desarrollado íntegramente en \emph{español}, con una interfaz diseñada bajo criterios de minimalismo y claridad pedagógica.

La elección del dominio bidimensional no es solo una restricción de alcance, sino una decisión pedagógica fundamentada: el continuo plano ocupa el punto intermedio idóneo entre el análisis unidimensional de barras y el análisis tridimensional de sólidos. El elemento de barra, con el que suele iniciarse la enseñanza del método, resulta insuficiente para ilustrar los conceptos centrales de la formulación isoparamétrica —el mapeo entre coordenadas naturales y físicas, la matriz Jacobiana, la matriz de deformación-desplazamiento y la cuadratura en dos direcciones—, que solo emergen al pasar al continuo. El análisis tridimensional, en el otro extremo, sí contiene esos conceptos, pero su aparato algebraico los oscurece: la matriz constitutiva crece a $6\times6$, la de deformación-desplazamiento a $6\times24$ en un hexaedro de ocho nodos y la rigidez elemental a $24\times24$ \autocite[p.~218]{cook2002concepts}. Esas dimensiones vuelven impracticable el seguimiento paso a paso que constituye el núcleo didáctico de la herramienta. El plano es la mínima dimensión en la que aparece el verdadero carácter de continuo y, a la vez, conserva matrices de tamaño legible —$\bm{D}$ de $3\times3$ y $\bm{B}$ de $3\times8$ en el elemento Q4—, lo que permite exhibir cada operación de forma trazable. Además, la formulación isoparamétrica, la cuadratura de Gauss y el ensamblaje son directamente generalizables a tres dimensiones \autocite{bathe2014fem,zienkiewicz2013fem}. El dominio plano actúa así como un puente: lo que el estudiante comprende en dos dimensiones se extiende de manera natural al caso tridimensional, planteado como línea de continuación del proyecto.

Las limitaciones del trabajo se reconocen de manera complementaria. El software no aborda análisis tridimensional, comportamiento no lineal (material o geométrico), efectos dinámicos ni dependientes del tiempo, ni otros tipos de elementos (triangulares, de orden superior a Q9, o elementos especiales). El catálogo de materiales se restringe a comportamiento elástico, isótropo y lineal, definido por el módulo de elasticidad y el coeficiente de Poisson. La malla se construye manualmente —por tablas, por dibujo en el lienzo o por importación DXF— o por generación estructurada en los casos de referencia: el software no incorpora mallado automático. La validación numérica se apoya en casos con solución analítica conocida y casos de referencia reconocidos de la literatura. La validación de la dimensión educativa se limita al contenido y a la coherencia del diseño, establecidos sobre la literatura y sobre un consenso publicado de expertos; la prueba de campo que mida el criterio de los estudiantes queda fuera del alcance temporal del trabajo y se deja diseñada en las recomendaciones. Estas fronteras responden a una decisión deliberada: privilegiar la transparencia y la solidez de un núcleo conceptual bien acotado por sobre la cobertura amplia de funcionalidades, coherente con el propósito didáctico de la herramienta.

\section*{Estructura del documento}

El presente documento se organiza en tres capítulos, además de esta introducción, de un capítulo final de conclusiones y de un conjunto de anexos. El \autoref{cap:marco_teorico} desarrolla el marco teórico: abre con los antecedentes y el estado del arte del software para la enseñanza del MEF —que sitúan el trabajo, justifican su oportunidad y cumplen el primer objetivo específico—, continúa con los fundamentos pedagógicos —los cuatro principios de aprendizaje que la exposición del canal de cálculo encarna— y desarrolla luego los fundamentos del método aplicados a la elasticidad lineal bidimensional (formulación variacional, funciones de forma, mapeo isoparamétrico, matriz Jacobiana, matriz de deformación, matriz constitutiva, cuadratura de Gauss, ensamblaje y recuperación de tensiones), la calidad de malla y los criterios de verificación y validación. El \autoref{cap:diseno} expone el diseño e implementación del modelo: abre con el diseño metodológico —el tipo de investigación y su modelo de simulación numérica, las variables y su matriz de consistencia, la población y los casos de estudio, el procedimiento con sus instrumentos, y los criterios de aceptación con que se contrasta la hipótesis— y continúa con el diseño e implementación de la herramienta —los requisitos y decisiones de diseño, el stack tecnológico, la arquitectura en torno al modelo de proyecto, el motor de cálculo, el pre-proceso interactivo, los módulos educativos, el post-proceso y la generación de la memoria de cálculo—. El \autoref{cap:resultados} presenta y analiza los resultados: expone el producto final, mide los datos, ejecuta la verificación y validación numérica contrastando EduFEM frente a soluciones analíticas y software comercial, establece la validez de contenido y la coherencia del diseño con la tabla de especificaciones y la matriz de principios, interpreta los hallazgos y contrasta la hipótesis. Los objetivos específicos se reparten entre los capítulos: el primero y el segundo se cumplen en el \autoref{cap:marco_teorico}; el tercero se materializa en el \autoref{cap:diseno}; y el cuarto y el quinto se cumplen en el \autoref{cap:resultados}, que reporta además la prueba de interoperabilidad. A continuación, el capítulo de conclusiones responde a cada objetivo en ese mismo orden, señalando el capítulo que lo desarrolla, discute el cumplimiento de la hipótesis y agrupa las recomendaciones por capítulo, con el diseño de la prueba de campo entre ellas. Sigue la lista de referencias bibliográficas. Las citas y esa lista siguen la norma Vancouver (ICMJE y NLM), con numeración arábiga por orden de primera mención —entre corchetes en lugar de superíndice, para poder acompañar el número con la página citada—, títulos de revista completos y conjunción antes del último autor; el localizador de página acompaña a toda cita que sostiene una ecuación, un valor numérico, un umbral, una definición o una atribución concreta, y se omite en las citas de encuadre. La presentación del documento sigue la séptima edición de las normas APA. Por último, los anexos reúnen el manual de instalación (\autoref{anexo:instalacion}), el manual de uso (\autoref{anexo:manual}), una selección de listados de código del motor numérico (\autoref{anexo:codigo}), los resultados extendidos de verificación y validación (\autoref{anexo:vyv}), el modelo de validación en SAP2000 y la solución analítica de Timoshenko (\autoref{anexo:sap2000}), los formatos de interoperabilidad (\autoref{anexo:formatos}) y una memoria de cálculo resuelta paso a paso sobre el ejemplo canónico (\autoref{anexo:memoria}).
